\documentclass[10pt,a4paper]{amsart}
\usepackage[margin=19mm]{geometry}
\usepackage{amsmath,amssymb,mathtools}
\usepackage[scr=rsfs,cal=euler]{mathalfa}
\usepackage{xfrac}
\usepackage[unicode,hidelinks,bookmarks=false]{hyperref}
\newcommand{\ie}{i.e.\ }
\newcommand{\cf}{cf.\ }
\newcommand{\resp}{respectively\ }
\newcommand{\Subsection}{\subsection}
\providecommand{\nobreakdash}{\nobreak\hbox{-}}
\setlength{\emergencystretch}{2em}
\multlinegap0pt
\usepackage{amssymb}
\newtheorem{thm}{Theorem}
\newcommand{\Pcal}{\mathcal{P}}
\newcommand{\gradb}{\nabla_{\Pcal}}
\title{Why the Kellogg Mesh Is Radial: A Mathematical Explanation of a Classical Computational Benchmark}
\author{Shun Zhang}
\date{Source: arXiv:2608.22991v2 (CC BY 4.0)}
\begin{document}
\maketitle

\noindent\emph{Source and licence.} Why the Kellogg Mesh Is Radial: A Mathematical Explanation of a Classical Computational Benchmark. Version arXiv:2608.22991v2 (CC BY 4.0), distributed by arXiv under the Creative Commons Attribution 4.0 International licence, http://creativecommons.org/licenses/by/4.0/, as recorded in the arXiv licence metadata for this exact version and shown on the abstract page https://arxiv.org/abs/2608.22991v2. Full licence notice: \texttt{licenses/arxiv-2608.22991-CC-BY-4.0.txt}.\par\smallskip

\noindent\textit{Editorial scope.} Counted unit: the article's thm thm:main (source0.tex lines 757--768). The article's complete original proof of it is delivered with it (source0.tex lines 770--776, one \texttt{proof} environment of 7 lines). No mathematics is written, repaired or re-derived: the statement and the proof are the article's own lines, in the article's own notation, and no retained line is substituted or re-flowed.  Retained uncounted as the source-local context the proof needs: lines 657--661; lines 668--671; lines 747--750.  Nothing was omitted from any retained span, so the package carries no omission record.  No retained line contains a citation, so the wrapper carries no bibliography and no bibliography entry.  No retained line contains a figure environment, an image-inclusion command or drawing code, so no image is delivered and the package declares no asset dependency.  Editorial formatting only, in addition: an AMS \texttt{amsart} preamble that restates the article's own declarations and macros used by the retained lines, namely the environments \texttt{thm} on separate counters, together with the standard packages those lines need.  The retained environments print in this document as Theorem 1 in order of appearance, whereas the article prints the same items with the numbering of its own full document.  The complete native source of the article as distributed in the arXiv e-print for this exact version is bundled unchanged as \texttt{sources/208382/source0.tex}.
\begin{equation}\label{eq:sector-gradient}
 |\nabla u|^2
 =r^{2\gamma-2}\bigl(\gamma^2\mu^2+(\mu')^2\bigr)
 =\gamma^2A_i^2r^{2\gamma-2}.
\end{equation}

\begin{equation}\label{eq:sector-hessian}
 |\nabla^2u|_F^2
 =2\gamma^2(1-\gamma)^2A_i^2r^{2\gamma-4}.
\end{equation}

\begin{equation}\label{eq:amplitude-balance}
 \boxed{
 \kappa_1A_1^2=\kappa_2A_2^2=\kappa_3A_3^2=\kappa_4A_4^2.}
\end{equation}

\begin{thm}[Hidden energy symmetry]\label{thm:main}
For Kellogg's perpendicular two-phase checkerboard and its homogeneous singular mode, there is a constant $\Lambda>0$ such that
\begin{equation}\label{eq:main-gradient}
  \boxed{\kappa(x)|\nabla u(x)|^2=\Lambda\, r^{2\gamma-2},}
\end{equation}
and
\begin{equation}\label{eq:main-hessian}
  \boxed{\kappa(x)|\gradb^2u(x)|_F^2
  =2(1-\gamma)^2\Lambda\, r^{2\gamma-4}.}
\end{equation}
Both right-hand sides are functions of $r$ alone.  Neither the angular variable nor the coefficient appears there, although each of $\kappa$ and $u$ depends on both.  The identities hold in the four open quadrants, and by the amplitude balance \eqref{eq:amplitude-balance} the one-sided traces of the two scalar quantities agree across every interface, so each admits a unique continuous extension to $\Omega\setminus\{0\}$.
\end{thm}

\begin{proof}
On $Q_i$, \eqref{eq:sector-gradient} and \eqref{eq:sector-hessian} give
\[
\kappa|\nabla u|^2=\kappa_i\gamma^2A_i^2\,r^{2\gamma-2}, \qquad \kappa|\gradb^2u|_F^2=2(1-\gamma)^2\bigl(\kappa_i\gamma^2A_i^2\bigr)r^{2\gamma-4}.
\]
By \eqref{eq:amplitude-balance} the quantity $\kappa_i\gamma^2A_i^2$ takes one common value on the four quadrants; call it $\Lambda$.
\end{proof}

\end{document}
